\section{Auxiliary Proofs and Exact Dependencies}
\label{app:algebra}
This appendix expands the mathematical ingredients used in the main text. Theorems~\ref{thm:main} and~\ref{thm:42} retain their original unrestricted code domains. E1, E2, and E3 are published inputs as specified in Section~2; Proposition~\ref{prop:computed} is the external finite computation.

\subsection{Puncturing, shortening, and parity}
For a set $S$ of $t$ coordinates, puncturing has kernel consisting of codewords supported in $S$. If $t<d(C)$, the kernel is zero, so dimension is preserved. Every remaining word loses at most $t$ ones, establishing the lower bound $d(C)-t$.

Let $\pi_S:C\to\F^S$ be coordinate projection, with rank $r$. The shortened code is the image of $\ker\pi_S$ after deleting $S$. This deletion is injective on the kernel because its members are zero on $S$. Rank-nullity gives dimension $\dim C-r$. No weight changes. For one coordinate, $r\le1$; if an exact smaller dimension is needed, choose a subspace of that dimension. This includes a coordinate that was identically zero before shortening.

Parity extension is an injective linear map because its first $n$ coordinates recover its input. Its weight is $\wt(x)$ for even-weight $x$ and $\wt(x)+1$ for odd-weight $x$. The minimum distance is therefore at least the least even integer greater than or equal to the original lower bound. Puncturing any coordinate of an even $[44,10,\ge18]$ code preserves dimension and gives distance at least 17. This proves the equivalence used at the start of Section~\ref{sec:reduction} without assuming that the punctured coordinate was originally a parity coordinate.

\subsection{A proof of the binary Griesmer bound}
For completeness, the needed inequality can be obtained by induction on dimension. Let $c\in C$ have minimum weight $D$, and puncture the $D$ positions in its support. The kernel is $\langle c\rangle$: any other nonzero word supported there has weight at least $D$ and thus equals $c$. The residual image has dimension $k-1$. For a nonzero residual word, take its two preimages $x,x+c$. If the common weight outside the support is $b$ and the weight of $x$ inside is $a$, then
\[
 a+b\ge D,\qquad D-a+b\ge D.
\]
Adding yields $b\ge\lceil D/2\rceil$. By induction,
\[
 n-D\ge\sum_{j=0}^{k-2}\left\lceil\frac{\lceil D/2\rceil}{2^j}\right\rceil
       =\sum_{j=1}^{k-1}\left\lceil\frac D{2^j}\right\rceil.
\]
The base $k=1$ is $n\ge D$. Monotonicity in $D$ gives \eqref{eq:griesmer} for a distance lower bound. In particular,
\[
 g_2(8,19)=19+10+5+3+2+1+1+1=42,
 \qquad g_2(7,20)=20+10+5+3+2+1+1=42.
\]
Both exceed 41 and 40, respectively. E1 plus the first inequality rules out distance at least 18 at parameters $(41,8)$.

\subsection{MacWilliams identity and the forcing polynomial}
For any $y\in\F^{44}$, character orthogonality gives
\[
 \sum_{x\in C}(-1)^{x\cdot y}=
 \begin{cases}2^{10},&y\in C^\perp,\\0,&y\notin C^\perp.\end{cases}
\]
In the second case, choose $x_0\in C$ with $x_0\cdot y=1$ and pair each $x$ with $x+x_0$. Summing over all $y$ of weight $j$, then interchanging two finite sums, yields
\[
 2^{10}B_j=\sum_{x\in C}\sum_{\wt(y)=j}(-1)^{x\cdot y}
 =\sum_{w=0}^{44} A_wK_j(w).
\]
For $z=44-2w$, the needed polynomials are
\[
 K_0(w)=1,\quad K_2(w)=\frac{z^2-44}{2},\quad
 K_4(w)=\frac{z^4-256z^2+5544}{24}.
\]
Substitution or direct expansion proves \eqref{eq:polynomial}. The potential negative intervals of $P(w)$ are $(18,20)$ and $(24,26)$, neither of which contains an even integer. Thus $P(w)\ge0$ at every allowed nonzero weight. Since $A_0=B_0=1$ and $B_2=0$,
\[
 \tfrac32\,2^{10}B_4+\tfrac{403}{2}\,2^{10}\ge224640,
 \qquad B_4\ge\frac{224640-206336}{1536}=\frac{143}{12}.
\]
This supplies the weight-four dual word without solving a linear program.

\subsection{Full support and the extension normal form}
The full-support claim uses the nondegenerate dot product on the message space $\F^{10}$. The shortened messages are $U^\perp$, where $U$ is the span of the four tail columns. A remaining column $v$ vanishes on all shortened messages if and only if $v\in(U^\perp)^\perp=U$. Because $\dim U=3$ and $v_4=v_1+v_2+v_3$, its nonzero vectors are exactly
\[
 v_1,v_2,v_3,v_4,\quad v_1+v_2,v_1+v_3,v_2+v_3.
\]
Each possibility contradicts independence of every three original columns, as explained in Lemma~\ref{lem:mother}. The shortening is therefore genuinely in the full-support class covered by E2.

The tail projection has kernel $K$ and image the dimension-three even subspace of $\F^4$. Choose lifts of $1001,0101,0011$ and prepend a basis for the kernel. This yields \eqref{eq:normal} by row operations, without restricting the code. For $s\in\F^3$ and $c\in K$, every word has the form
\[
 \left(c+s_1a_1+s_2a_2+s_3a_3,\ s_1(1001)+s_2(0101)+s_3(0011)\right).
\]
For $s=0$ the weight is at least 18 unless the word is zero. For $\wt(s)=1$ or 2 the tail weight is 2, and for $s=(1,1,1)$ it is 4. This proves necessity and sufficiency of \eqref{eq:extension}. Replacing an $a_i$ by a vector in its coset amounts to adding kernel rows, so all conditions depend only on cosets. A coordinate permutation of $K$ extends by the identity on the tail and preserves every condition.

\subsection{Unique even representatives and pair invariance}
In systematic form, $(r_i)_{p_j}=\delta_{ij}$. Hence $N_K(x)$ in \eqref{eq:normalizer} is zero at every pivot and $N_K(x+r_i)=N_K(x)$. Conversely, $N_K(x)=0$ implies $x=\sum_i x_{p_i}r_i\in K$. The map induces a bijection between the quotient and its zero-pivot image. Since the rows of $K$ are even, it preserves parity. On the zero-pivot subspace, evenness is the nonzero equation that the 33 free bits sum to zero; its dimension is 32. All coset and pair distances are unchanged under this normalization because $N_K(x)+x\in K$ and the map is linear.

\subsection{Distance normalization and the fifth target}
The condition $d\ge1$ in Lemma~\ref{lem:normalize} is essential to its dimension argument. If one punctured all $D$ positions of a minimum word to obtain target zero, that word would enter the kernel. For positive target $d$, only $D-d<D$ positions are removed, so no nonzero word is lost. Zero padding is permitted because our original domain is the full class of binary linear codes.

After parity extension of a hypothetical $[42,16,\ge13]$ code, two one-coordinate shortenings give dimension at least 14 and distance at least 14. Selecting a 14-dimensional subspace preserves that distance lower bound. Normalization then produces exact distance 14 at the same length and dimension. This argument discharges the exact-distance hypothesis of E3; no assertion about preserving a chosen minimum word under shortening is needed.

\section{Certificate Semantics and Reproduction}
\label{app:reproduction}
\subsection{Finite checks and completeness}
A complete checker for a mother $K$ must verify three independent properties:
\begin{enumerate}
\item The supplied matrix defines an even, full-support $[40,7,18]$ code and matches the coordinates used by the scan.
\item The list $F_K$ equals the entire set of even cosets at distance at least 16; membership and completeness are separate obligations.
\item Every unordered pair of distinct entries has a message satisfying \eqref{eq:witness}.
\end{enumerate}
The two scans provide the second property with different systematic bases and enumeration orders. The pair checker reads exactly $\binom{|F_K|}{2}$ witnesses in the declared pair order. Its weight histogram, over all mothers, is
\[
\begin{array}{c|rrrrrrr}
\text{rejection weight}&2&4&6&8&10&12&14\\\hline
\text{count}&186172&922339&4503729&16924956&61595955&171053508&363867130.
\end{array}
\]
These counts sum to 619,053,789. A message value outside $0,\ldots,127$, a missing or trailing record, a wrong matrix binding, or weight at least 16 is a verification failure.

\subsection{What is supplied with this draft}
The accompanying data directory supplies the literal mother matrices, complete per-case coverage metadata, historical pair-replay summary, and lower generators. A dependency-free script recomputes all lower spectra, mother parameters, the forcing-polynomial identity at every integer weight, and coverage counts. A separate catalogue check uses graph canonicalization for pairwise inequivalence. The local evidence map identifies the existing full scan outputs and the approximately 1.68 GB of pair-witness streams; these large streams are not embedded in the compact manuscript-source archive. Historical absolute paths in metadata are resolved against the current project root, with original records preserved.

During manuscript preparation, we rebuilt both Lean projects and reran their audits, re-enumerated every lower-bound matrix, rechecked all mother parameters and pairwise inequivalence, compared both saved accepted sets with their coordinate bindings, and replayed all 619,053,789 pair witnesses. All these checks passed. The validation log distinguishes these new checks from the original full-domain scans. Rebuilding the Lean projects or rechecking matrix spectra does not re-execute the $172\cdot2^{32}$ coset domains. Replaying all saved pair witnesses likewise does not certify that a saved list is exhaustive. Reproduction commands make these scopes explicit.

\section{Formal Interfaces and Imported Facts}
\label{app:formal}
The first project exposes the following theorem. Names are reproduced verbatim; line breaks below are typographical.
\begin{lstlisting}[basicstyle=\ttfamily\small,breaklines=true]
ECC.mainClaim_of_external_results
  (hE1 : PublishedNonexistence41)
  (hE2 : PublishedClassBound)
  (hDistinct : CatalogueInequivalent)
  (hComputed : forall i : Fin 172, PairIncompatible (catalogueCode i))
  : MainClaim
\end{lstlisting}
Here \texttt{MainClaim} means nonexistence of an actual dimension-ten subspace of $\F^{43}$ with every nonzero word of weight at least 17. E1 excludes exact $[41,8,18]$; E2 bounds the number of full-support exact $[40,7,18]$ classes. The theorem proves the connection from these hypotheses to arbitrary target codes, including the full-support reduction and actual matrix catalogue. No unproved ``catalogue reduction'' parameter remains in this final interface.

The corollary interface is \texttt{ECC.four\_table\_consequences\_of\_external\_results}. Its conclusion combines the three exact entries for $(43,10),(44,10),(44,11)$ with existence of $[45,11,\ge16]$ and nonexistence of $[45,11,\ge18]$. Kernel checks of packed matrices use \texttt{decide +kernel}. The exhaustive coset scans and graph canonicalization are external to these kernel computations.

The second project exposes
\begin{lstlisting}[basicstyle=\ttfamily\small,breaklines=true]
ECC42.reduction_42_16_to_41_14
  : CodeExists 42 16 13 -> ExactCodeExists 41 14 14
ECC42.nonexistence_of_proposition2
  : Proposition2 -> not (CodeExists 42 16 13)
\end{lstlisting}
where \texttt{Proposition2} is the published exclusion of exact $[41,14,14]$. The imported proposition is an explicit argument, not an axiom declaration. Both developments use Lean 4.34.1; their pinned manifests are included with the corresponding local projects. Audit commands inspect both final theorem types and transitive axioms, since an axiom list alone cannot reveal outstanding theorem parameters.

\section{Source Provenance and AI Attribution}
The published inputs are \citet{bouyukliev2000}, \citet{bouyukliev2001}, and Proposition~2 (printed page 179) of \citet{bouyuklieva2020}. The archived source assessment for E1 additionally locates Theorem~6.3 in Bouyukliev's doctoral thesis, and for E2 locates Table~6.2. These thesis numbering references are not journal theorem numbers. The 2025 hull paper provides a separate count cross-check. Historical classification internals are not presented as rerun computations.

The provenance records for the AI interventions retain original responses, subsequent actions, and reconciliation notes. The early Gemini profiling suggestion is connected to the native-filter experiment; the Grok completeness request is connected to the second full scan and set comparison. Later reviews are recorded as reviews of already assembled evidence. Model agreement is not included among the hypotheses of either mathematical theorem. The local record permits these specific attributions, but does not support a percentage allocation of the proof to humans or models, or a controlled comparison among systems.
